%==============================================================================
% A Glass Box Around the Black Box:
% Deterministic Auditing of LLM Answers Without Inference
%
% Karthik Barma, Sarita Singh.  DRAFT for author revision.
% Evidence point: commit d852db090cd0b72a50b1738afca9392332445810
%
% Self-contained two-column US Letter layout: compiles on a plain TeX Live
% install with no venue style file. To retarget an AAAI venue, drop in the
% author kit, replace the preamble down to \begin{document}, keep the body.
%==============================================================================
\documentclass[10pt,twocolumn,letterpaper]{article}

\usepackage[letterpaper,margin=0.75in,columnsep=0.30in]{geometry}
\usepackage[T1]{fontenc}
\usepackage[utf8]{inputenc}
\usepackage{times}
\usepackage{microtype}
\usepackage{amsmath,amssymb}
\usepackage{graphicx}
\usepackage{booktabs}
\usepackage{array}
\usepackage{xcolor}
% Ragged-right fixed-width column. Justified text in narrow table cells is what
% produces the stretched interword gaps; never justify inside a p-column here.
\newcolumntype{L}[1]{>{\raggedright\arraybackslash}p{#1}}
\usepackage[font=small,labelfont=bf,skip=5pt]{caption}
\usepackage{enumitem}
\usepackage{tikz}
\usetikzlibrary{shapes.geometric,arrows.meta,positioning,matrix,fit,calc}
\usepackage[raggedright]{titlesec}
\usepackage[hidelinks,breaklinks]{hyperref}
\usepackage{url}

%------------------------------------------------------------------ typography
% Headings: ragged right and never hyphenated. Long section titles were
% breaking as "Five Lev-els" and "Case Stud-ies" in the narrow columns.
\titleformat{\section}{\normalfont\large\bfseries\raggedright}{\thesection}{0.65em}{}
\titleformat{\subsection}{\normalfont\normalsize\bfseries\raggedright}{\thesubsection}{0.6em}{}
\titlespacing*{\section}{0pt}{2.4ex plus 1ex minus .2ex}{1.3ex plus .2ex}
\titlespacing*{\subsection}{0pt}{2.0ex plus 1ex minus .2ex}{1.0ex plus .2ex}
\pretolerance=2000
\tolerance=3000
\emergencystretch=2em
\hbadness=2000

% The twocolumn class option turns on \flushbottom, which forces both columns to
% the same height by stretching the glue between paragraphs. On a page carrying a
% tall float that stretch collects into a few large visible gaps. Ragged bottom
% keeps paragraph spacing uniform and lets column depth vary instead.
\raggedbottom
% Run-in \paragraph headings default to 3.25ex of stretchable space above; fix it
% so those headings sit at a constant distance from the preceding text.
\titlespacing*{\paragraph}{0pt}{1.7ex plus 0.3ex minus 0.2ex}{0.9em}
\setlength{\parskip}{0pt plus 0.4pt}
\setlength{\floatsep}{9pt plus 2pt minus 2pt}
\setlength{\textfloatsep}{11pt plus 2pt minus 3pt}
\setlength{\intextsep}{9pt plus 2pt minus 2pt}
\setlength{\dblfloatsep}{9pt plus 2pt minus 2pt}
\setlength{\dbltextfloatsep}{11pt plus 2pt minus 3pt}

\definecolor{gbInk}{HTML}{1B2A33}
\definecolor{gbBox}{HTML}{EEF3F6}
\definecolor{gbEdge}{HTML}{5A7684}
\definecolor{gbOk}{HTML}{2F7D32}
\definecolor{gbOkBg}{HTML}{DFF5E1}
\definecolor{gbPend}{HTML}{A17B12}
\definecolor{gbPendBg}{HTML}{FDF3D8}
\definecolor{gbFut}{HTML}{6B6B6B}
\definecolor{gbFutBg}{HTML}{EFEFEF}

\newcommand{\lv}[1]{\mbox{\textbf{L#1}}}
\newcommand{\code}[1]{\mbox{\texttt{\small #1}}}
\newcommand{\probe}[1]{\mbox{\texttt{\scriptsize #1}}}
\newcommand{\rt}[1]{\mbox{\texttt{\scriptsize #1}}}
\hyphenation{Glass-Box de-ter-min-is-tic mar-ket-place}

\title{\bfseries A Glass Box Around the Black Box:\\
Deterministic Auditing of LLM Answers Without Inference}

\author{
\begin{tabular}{c@{\hspace{2.4em}}c}
\textbf{Karthik Barma}\thanks{First author. GlassBox is developed under \emph{Aura}, an
unregistered umbrella brand and not a company or legal entity. The system charges no
fee, accepts no payments, and sells nothing.} &
\textbf{Sarita Singh} \\
\small MS Artificial Intelligence & \small Associate Teaching Professor \\
\small Khoury College of Computer Sciences & \small Khoury College of Computer Sciences \\
\small Northeastern University & \small Northeastern University \\
\small \texttt{barma.k@northeastern.edu} & \small \texttt{s.singh@northeastern.edu} \\
\end{tabular}
}
\date{}

%==============================================================================
\begin{document}
\maketitle

\begin{abstract}
\noindent
Tools that audit AI answers generally assume a model in the loop: a judge model, a
classifier, or a guardrail service carrying its own inference cost. That assumption makes
auditing something a deployment enables deliberately, for a subset of traffic, rather
than something that can be on by default. We ask how much useful auditing survives when
the auditor is permitted no model inference at all.

We present GlassBox Lite, a deterministic structural auditor that examines an explicitly
submitted question and answer pair using seven bounded probes: claim segmentation,
allowlisted arithmetic recomputation, direct lexical contradiction, unsupported absolute
certainty, citation transparency, external fact-check scope, and prompt-injection
signals. It issues no model call and performs no network lookup. It never inspects the
generating model, so it places a glass box \emph{around} a black box rather than opening
one, and it emits a compact Trust Card carrying a verdict, a structural score, per-probe
outcomes, and explicit limitations.

Removing inference has a consequence beyond cost. Because identical inputs yield
identical outputs, determinism becomes a \emph{deployment} property as much as an
evaluation one, making idempotency, replay suppression, and reproducible audit records
exact rather than approximate. We confirmed this against the live public endpoint, where
three identical production requests returned byte-identical responses.

We then operate this auditor behind eight heterogeneous surfaces and report what
deployment actually costs. The public Trust Card projection withholds submitted content
and every verifier-internal field, verified by planting sentinel tokens in a live request
and confirming their absence. All provider routes fail closed on unsigned requests. At
the evaluated commit, 108 automated tests pass across six suites and twelve
continuous-integration jobs succeed on the merged main commit.

Our most transferable contribution is methodological. Deployment claims in this area
routinely collapse five distinct states into the word ``available.'' We separate
implementation presence, provider registration, public installability, production
end-to-end evidence, and directory approval into an explicit five-level scale, and report
our own results against it, including the surfaces that sit low on it.

We then measure the probes themselves on a constructed benchmark of 187 items with labels
fixed by construction. Precision is 1.000 on every probe with zero false positives, but
recall splits the probes in two: those that \emph{compute} a property generalise to unseen
phrasing (arithmetic recomputation 1.000, contradiction 0.833) while those that \emph{match
a vocabulary} do not (0.000 to 0.333 on a held-out split). Repairing the vocabularies raised
development recall to 1.000 and moved held-out recall almost not at all. That is a bound on
the zero-inference design, not a defect list. GlassBox remains a structural auditor and
never a fact-checker: we measure whether probes fire correctly, not whether answers are
true.
\end{abstract}

%==============================================================================
\section{Introduction}

An answer auditor that costs money per audit is an auditor that gets switched off. If
verification requires a paid model API, the verification budget scales with traffic, and
the rational deployment response is to sample: audit some answers, some of the time, in
some contexts. The cases most worth auditing then become the ones most likely to be
missed.

This paper explores the opposite corner of the design space. We constrain the auditor to
perform no model inference whatsoever and ask what remains. The honest answer is: less
than a model-based judge~\cite{zheng2023judging}, but more than nothing, and what remains
has properties a model-based judge cannot have.

What remains is \emph{structural} auditing. GlassBox Lite does not decide whether an
answer is true. It decides whether an answer is internally well formed: whether its
arithmetic recomputes, whether it contradicts itself, whether it asserts certainty it has
not earned, whether it cites anything, and whether it contains instruction-like text that
a downstream client should treat as inert. These properties are narrow. They are also
checkable without a model, and checkable identically every time.

The metaphor in our title is deliberate and load-bearing. Interpretability research
opens the black box, and glass-box modelling replaces it with an inherently transparent
one. We do neither. GlassBox never inspects weights, activations, or model internals, and
it does not require the model to explain itself, which is fortunate given the evidence
that such explanations can be unfaithful~\cite{turpin2023unfaithful} and that model
self-report is a fragile audit substrate~\cite{kadavath2022know}. It places a transparent
enclosure around an opaque generator: transparency about the \emph{answer}, with zero
access to the \emph{model}.

That constraint is where the paper turns. A deterministic auditor is not merely cheaper
than a sampled one, it is operationally different. Retry handling becomes exact rather
than probabilistic. A duplicate provider delivery can be suppressed with confidence that
the suppressed result would have been identical. An audit record can be reproduced from
its inputs. None of this is available to a service whose verdict is drawn from a sampling
distribution, which is the position of sampling-based detectors such as
SelfCheckGPT~\cite{manakul2023selfcheckgpt}.

We therefore treat deployment as the primary object of study rather than an afterthought.
GlassBox runs as a single verification service behind a web application, a remote Model
Context Protocol (MCP) endpoint~\cite{mcp2025spec}, a GitHub App, Discord, Telegram,
Slack, an authenticated API, a browser extension, two IDE clients, a Notion connection,
and a Reddit Devvit app. Building the verifier took a fraction of the effort. Satisfying
eight providers' authentication schemes, invocation models, consent requirements,
acknowledgement deadlines, and review processes took the rest.

That asymmetry produced our second contribution, a reporting discipline. Partway through
this work it became clear that ``GlassBox is available on Discord'' and ``GlassBox is
available on GitHub'' were both true sentences describing very different situations, and
that no vocabulary distinguished them. We introduce a five-level evidence scale
(Section~\ref{sec:method}) and hold every claim to it, including the unflattering ones.

\paragraph{Contributions.}
\begin{enumerate}[leftmargin=1.15em,itemsep=1pt,topsep=2pt]
\item A deterministic, zero-inference structural auditor with seven bounded, explainable
probes, demonstrated byte-identical on a live public endpoint.
\item A privacy-minimized public Trust Card projection that withholds submitted content
and all verifier-internal metadata, verified by planted-token testing against production.
\item A cross-platform gateway preserving provider-native authentication, explicit
invocation, tenant admission, idempotency, rate limits, deadlines, and output
neutralization, with every provider route verified to fail closed.
\item A probe-level benchmark with a held-out split, yielding the finding that model-free
probes generalise exactly when they compute a property and barely at all when they match a
lexicon.
\item Reproducibility evidence at a fixed public commit, and a five-level evidence
discipline under which negative results are reported as first-class findings.
\end{enumerate}

\paragraph{Non-goals.}
GlassBox is not a fact-checker, a source authenticator, a truth oracle, a hallucination
detector with measured accuracy, a semantic contradiction engine, a general-purpose
calculator, or a complete prompt-injection security boundary. The precision and recall
reported in Section~\ref{sec:bench} describe whether each probe fires when it should on a
constructed benchmark; they say nothing about how often these failure modes occur in real
traffic, and nothing about factual truth. Section~\ref{sec:limits} states these limits in
full.

%==============================================================================
\section{Problem Definition}

Consider an answer containing the expression $9 \times 9 = 80$. Detecting this requires
no world knowledge, no retrieval, and no model. It requires parsing an arithmetic
expression and recomputing it. Yet in most deployed architectures, catching it means
either a second model call or nothing at all.

Structural errors of this kind form a real and under-served class. An answer may
contradict itself between its second and fifth sentences. It may assert that something is
``certainly'' true while supplying no support. It may reference sources that do not exist
as citations at all. It may contain text engineered to be read as an instruction by
whatever consumes it~\cite{greshake2023injection,perez2022ignore}. None of these require
establishing external truth, and all are visible in the structure of the answer itself.
This is a deliberately narrower target than the hallucination problem as usually
framed~\cite{ji2023survey,huang2023hallucination}, and the narrowing is what makes a
model-free treatment possible.

We define the problem accordingly. \emph{Given an explicitly submitted question and
answer pair, produce a bounded, deterministic, explainable structural audit without model
inference, and deliver it across heterogeneous platforms without leaking the submitted
content.}

Three constraints shape everything that follows. \textbf{Zero paid API:} the verifier may
not call a paid model, which is what makes default-on auditing economically possible and
what forces the probes to be conservative. \textbf{Opt-in only:} GlassBox never audits
ambient traffic, it runs when a user explicitly invokes it. \textbf{Privacy-minimized
output:} on a public interface the result must not echo the submitted content nor expose
verifier internals, because the audit is useful precisely to the extent it can be shown
to someone without re-disclosing what was audited.

%==============================================================================
\section{Threat Model}

We consider five adversaries. A \textbf{hostile submitter} supplies content designed to
manipulate the auditor or the downstream client: prompt injection, hostile markup,
control and bidirectional characters, oversized payloads. A \textbf{forger} sends
unsigned or wrongly signed provider webhooks. A \textbf{replay adversary} resends valid
provider deliveries to duplicate work. A \textbf{resource adversary} floods the service to
exhaust a free-tier budget. A \textbf{curious observer} reads a public Trust Card hoping
to recover the submitted content or the verifier's internal reasoning.

Controls appear in Section~\ref{sec:controls} and their verification in
Section~\ref{sec:testing}. These are implementation-tested controls, not a security
assessment. We have conducted no penetration test, no static analysis at scale, and no
third-party audit, and we claim no compliance certification.

Out of scope: a malicious platform provider, a compromised host, network adversaries
beyond TLS, and any claim about the correctness of the audited answer's external facts.

%==============================================================================
\section{Trust Cards and the Structural Score}

A Trust Card is the unit of output: a verdict, a structural score, a claim count, a
finding count, a highest-severity marker, per-probe outcomes, and a caveat list.

The score is an Epistemic Confidence Score (ECS), and its name is the most dangerous
thing about it. \textbf{ECS is a structural score. It is not a probability that the answer
is true, and it is not calibrated against any labelled outcome.} An answer can score
highly and be entirely false, because a fluent, self-consistent, well-cited,
appropriately hedged falsehood passes every structural probe we run. We surface the
score's decomposition so a reader can see which dimensions drove it, and we attach
caveats to every card stating the limits of the audit within the card itself.

The verdict policy is deliberately conservative, on the view that over-trusting is the
more costly error. A single high-severity structural failure suffices to reject
regardless of the aggregate score. Our canonical canary illustrates this: the answer
$9 \times 9 = 80$ yields $\mathrm{ECS} = 0.8143$, high, because the answer is short,
self-consistent, appropriately scoped, and free of injection signals, while the verdict is
\code{reject}, driven by one high-severity \probe{arithmetic\_sanity} failure. The score
describes structure, the verdict describes acceptability, and they are not the same
question.

%==============================================================================
\section{The Lite Verifier}
\label{sec:lite}

Lite runs seven probes over a bounded input of 6{,}000 characters for the question,
12{,}000 for the answer, and at most eight optional intents (Table~\ref{tab:probes}).

\begin{table}[t]
\centering
\small
\caption{The seven probes. All are model-free and network-free, and all are conservative:
where a probe cannot reach a confident structural judgement, it passes. The deliberate
limitation of each probe is given in italics. The system under-reports rather than
over-reports, because an auditor that cries wolf gets disabled, and a disabled auditor has
precision zero.}
\label{tab:probes}
\setlength{\tabcolsep}{4pt}
\renewcommand{\arraystretch}{1.2}
\begin{tabular}{@{}L{0.405\columnwidth}L{0.535\columnwidth}@{}}
\toprule
\textbf{Probe} & \textbf{What it checks, and what it cannot} \\
\midrule
\probe{claim\_extraction}       & Atomic claims within a fixed limit.
                                  \emph{Cost bounded by design, not by answer length.} \\
\probe{arithmetic\_sanity}      & Recomputes allowlisted expressions.
                                  \emph{Silent outside the allowlist.} \\
\probe{internal\_contradiction} & Direct lexical and repeated-value conflicts.
                                  \emph{Not inference-grade.} \\
\probe{unsupported\_certainty}  & Certainty language without support.
                                  \emph{A lexical signal only.} \\
\probe{citation\_verifiability} & Citation presence and transparency.
                                  \emph{Citations are never authenticated.} \\
\probe{fact\_check\_scope}      & Answers overstating what a non-web audit establishes.
                                  \emph{Scope only.} \\
\probe{prompt\_injection}       & Instruction-like text, held inert.
                                  \emph{A signal, never containment.} \\
\bottomrule
\end{tabular}
\end{table}

Two design choices deserve comment. First, \probe{arithmetic\_sanity} uses an allowlist
rather than a general expression evaluator, because a general evaluator over untrusted
input is both a correctness and a security liability, and because a false negative is
cheaper here than a false positive. Second, \probe{prompt\_injection} treats matched text
strictly as data to be reported on, never as instruction to be followed: an auditor that
executed the content it audits would be the attack surface rather than the defence.

\paragraph{Determinism.}
No probe calls a model or the network, so identical inputs yield identical outputs. We
verified this against production rather than asserting it from source. Three identical
requests to the live public endpoint returned byte-identical response bodies, with SHA-256
digest prefix \code{f33f8274}. We note the contrast with our own earlier six-tool MCP
server, which is model-assisted and for which we make no determinism claim beyond
canonical assembly and hashing (Section~\ref{sec:crosslang}).

%==============================================================================
\section{The Public Projection}
\label{sec:privacy}

The public MCP projection is a deliberate subtraction. It omits the question, the answer,
extracted excerpts, verifier evidence, free-form rationale, timestamps, audit and log
identifiers, input hashes, and request or session metadata. What survives is the verdict,
the score, counts, fixed-category findings, per-probe outcomes, and caveats: enough to act
on, not enough to reconstruct the input.

We verified this behaviourally rather than by inspection. We issued a live production
request with unique sentinel tokens embedded in both the question and the answer, and
confirmed that neither token appeared anywhere in the response, and that none of
\code{log\_id}, \code{inputs\_hash}, \code{generated\_at}, \code{timestamp}, or
\code{excerpt} was present. We recommend this test generally: projection policies are easy
to state, easy to implement correctly once, and easy to regress silently when a field is
added upstream.

\paragraph{What we do not claim.}
Platform providers and the infrastructure host receive and process request traffic. We
therefore make no claim of ``zero data collection,'' ``nothing leaves the device,'' ``no
third parties,'' or end-to-end encryption, and we regard such claims in comparable systems
as generally unsupportable. Our retention posture is bounded, not absent: raw content is
not persisted to a database or to files by the gateway; content and results remain
transiently in memory for up to five minutes while work and delivery complete; completed
event identifiers are retained up to 24 hours for replay protection; and rate-control
identifiers persist up to ten minutes.

%==============================================================================
\section{Adapter Architecture}

One verification core sits behind a uniform pipeline (Figure~\ref{fig:arch}). Adapters
differ only at the ends: how a provider authenticates, and how a result is rendered.
Everything between is shared, which is what makes eight surfaces maintainable by one
developer.

\begin{figure}[t]
\centering
\begin{tikzpicture}[
  font=\scriptsize,
  stage/.style={draw=gbEdge,fill=gbBox,rounded corners=1.8pt,
                text width=0.82\columnwidth,align=center,inner sep=3.4pt,text=gbInk},
  client/.style={draw=gbEdge,fill=white,rounded corners=1.4pt,inner sep=0pt,
                 text=gbInk,minimum height=4.6mm,minimum width=15.5mm,align=center},
  ar/.style={-{Stealth[length=1.6mm]},draw=gbEdge,line width=0.45pt}
]
% Clients laid out in a matrix so both rows are centred on the same axis.
\matrix (cl) [matrix of nodes, nodes={client}, column sep=1.5mm, row sep=1.5mm,
              ampersand replacement=\&] {
  Web/PWA \& MCP \& GitHub \& Discord \\
  Telegram \& Slack$^{\dagger}$ \& API \& Devvit$^{\ddagger}$ \\
};

\node[stage,below=5.2mm of cl] (auth)
  {Provider authentication $+$ strict input bounds\\
   {\tiny Ed25519 / HMAC / secret header / bearer; unsigned $\rightarrow$ 401}};
\node[stage,below=3.4mm of auth] (admit)
  {Tenant admission $+$ requester limits $+$ replay protection\\
   {\tiny 10 audits / 10 min; completed event ids held 24\,h}};
\node[stage,below=3.4mm of admit] (queue)
  {Bounded queue, concurrency 1, daily ceiling\\
   {\tiny 100 accepted audits per day}};
\node[stage,below=3.4mm of queue,fill=gbOkBg,draw=gbOk,line width=0.7pt] (lite)
  {\textbf{GlassBox Lite deterministic verifier}\\
   {\tiny no model call, no network lookup, seven bounded probes}};
\node[stage,below=3.4mm of lite] (proj)
  {Privacy-minimized projection\\
   {\tiny drops question, answer, excerpts, timestamps, log\_id, inputs\_hash}};
\node[stage,below=3.4mm of proj] (fmt)
  {Platform-native Trust Card formatting\\
   {\tiny markdown, mentions, links, control and bidi characters neutralized}};

\draw[ar] (cl.south) -- (auth.north);
\draw[ar] (auth) -- (admit);
\draw[ar] (admit) -- (queue);
\draw[ar] (queue) -- (lite);
\draw[ar] (lite) -- (proj);
\draw[ar] (proj) -- (fmt);
\end{tikzpicture}
\caption{The GlassBox request path. A single deterministic verifier sits behind a uniform
admission, bounding, and projection pipeline; adapters vary only in authentication and
rendering. $^{\dagger}$Slack is an allowlisted single-workspace pilot.
$^{\ddagger}$Reddit Devvit is blocked pending outbound fetch-domain approval
(Section~\ref{sec:negative}).}
\label{fig:arch}
\end{figure}

The design tension is that providers disagree about nearly everything. Discord signs with
Ed25519 and expires interaction tokens on a deadline. Slack and GitHub sign with HMAC.
Telegram uses a secret header. MCP clients speak JSON-RPC over Streamable
HTTP~\cite{mcp2025spec}. Reddit Devvit requires an approved outbound-fetch allowlist. The
gateway normalises these into a single admission decision, and stops delivery before a
provider's token expires rather than emitting a late result.

%==============================================================================
\section{Authentication and Abuse Controls}
\label{sec:controls}

Every provider route validates provider-native authentication against the preserved raw
request body and \emph{fails closed}. Wrong events, missing delivery identifiers, unsigned
requests, and forged requests are rejected before any verification work is queued.

Admission is explicit. Public availability is per-platform rather than global: the deployed
configuration opens a selected subset and leaves the remainder allowlisted. Rate limits run
\emph{before} verifier invocation, at ten audits per ten minutes per requester, under a
global ceiling of 100 accepted audits per day, with verifier concurrency pinned at one.
Completed event identifiers are reserved through delivery and released on delivery failure,
so a provider retry after a genuine failure still succeeds while an in-flight retry does not
produce a duplicate reply.

Credential handling is minimal by construction. The MCP child process receives an explicit
environment allowlist rather than platform credentials; marketplace payloads and OAuth
tokens are neither logged nor persisted; a temporary OAuth token is revoked immediately
after use; and OAuth state is signed, one-time, plan-bound, and expires after ten minutes.

Output is neutralized before delivery: markdown, mentions, links, control characters, and
bidirectional characters are defanged, Discord mentions are disabled, and Slack link
unfurling is off. An auditor that rendered hostile submitted content into a privileged
channel would itself become the attack.

%==============================================================================
\section{Five Levels of Evidence}
\label{sec:method}

We found no adequate vocabulary for what we needed to report, so we adopted one
(Table~\ref{tab:levels}). The alternative is a single word, ``tested,'' spanning an
enormous range of epistemic strength: a system with a green test suite and a system with a
publicly viewable production result are not in the same evidentiary position.

\begin{table}[t]
\centering
\small
\caption{The evidence scale used throughout this paper. Every claim we make is tagged with
one of these levels, and Figure~\ref{fig:ladder} places each deployed surface on it.}
\label{tab:levels}
\setlength{\tabcolsep}{5pt}
\begin{tabular}{@{}cL{0.79\columnwidth}@{}}
\toprule
\textbf{Level} & \textbf{Meaning} \\
\midrule
\lv{4} & Public production end-to-end: a real request traversed a public surface and
produced a captured, publicly viewable result. \\
\lv{3} & Production signed or configured canary: an operator-initiated request against the
live production route, correctly authenticated or deliberately forged. \\
\lv{2} & Automated test, local or continuous integration, including adversarial tests. \\
\lv{1} & Implementation and configuration present, not exercised end to end. \\
\lv{0} & Proposed or future work. \\
\bottomrule
\end{tabular}
\end{table}

We answer no accuracy question at any level, because that requires a labelled benchmark we
have not built (Section~\ref{sec:future}).

\begin{figure}[t]
\centering
\begin{tikzpicture}[
  font=\scriptsize,
  lvl/.style={draw=gbEdge,rounded corners=1.4pt,minimum width=7.5mm,
              minimum height=5.2mm,align=center,text=gbInk},
  ev/.style={rounded corners=1.4pt,text width=0.655\columnwidth,align=left,
             inner sep=3.2pt,text=gbInk,draw=gbEdge}
]
\node[lvl,fill=gbOkBg,draw=gbOk] (l4) {\lv{4}};
\node[ev,fill=gbOkBg,draw=gbOk,right=2.0mm of l4]
  {Web/PWA; remote MCP (discovery, execution, determinism, projection); GitHub App};

\node[lvl,fill=gbPendBg,draw=gbPend,below=2.0mm of l4] (l3) {\lv{3}};
\node[ev,fill=gbPendBg,draw=gbPend,right=2.0mm of l3]
  {Discord; Telegram; authenticated API; GitHub Marketplace (submitted)};

\node[lvl,fill=gbPendBg,draw=gbPend,below=2.0mm of l3] (l2) {\lv{2}};
\node[ev,fill=gbPendBg,draw=gbPend,right=2.0mm of l2]
  {Browser extension; Notion; VS Code; JetBrains; Devvit; Slack pilot};

\node[lvl,fill=gbFutBg,draw=gbFut,below=2.0mm of l2] (l1) {\lv{1}};
\node[ev,fill=gbFutBg,draw=gbFut,right=2.0mm of l1]
  {ChatGPT directory; Claude directory; npm six-tool MCP};

\node[lvl,fill=gbFutBg,draw=gbFut,below=2.0mm of l1] (l0) {\lv{0}};
\node[ev,fill=gbFutBg,draw=gbFut,right=2.0mm of l0]
  {Chrome Web Store; Mattermost, Discourse, Teams; labelled benchmark};
\end{tikzpicture}
\caption{GlassBox surfaces on the evidence scale at the evaluated commit. Reporting the
lower rows is the point: a paper that promotes every row to ``available'' destroys the
information a reader needs in order to judge it.}
\label{fig:ladder}
\end{figure}

%==============================================================================
\section{Adversarial Testing}
\label{sec:testing}

At the evaluated commit, 108 automated tests pass across six heterogeneous suites
(Table~\ref{tab:tests}). This is the platform-layer total and not the repository's whole
test count.

\begin{table}[t]
\centering
\small
\caption{Test coverage at the evaluated commit. All suites except JetBrains were
re-executed locally during preparation of this paper; JetBrains is a Gradle project and
rests on continuous-integration evidence.}
\label{tab:tests}
\begin{tabular}{@{}lrl@{}}
\toprule
\textbf{Suite} & \textbf{Tests} & \textbf{Evidence} \\
\midrule
Gateway                & 69 & re-run locally \\
Notion (native)        & 12 & re-run locally \\
Browser extension      & 10 & re-run locally \\
Reddit Devvit          & \phantom{0}9 & re-run locally \\
VS Code                & \phantom{0}5 & re-run locally \\
JetBrains              & \phantom{0}3 & CI job green \\
\midrule
\textbf{Platform total} & \textbf{108} & \\
\bottomrule
\end{tabular}
\end{table}

The gateway suite is weighted toward adversarial and operational behaviour rather than
happy paths: forged and unsigned requests, wrong events, missing delivery identifiers,
in-flight retry duplication, rate limits ahead of invocation, the global ceiling, stalled
verification with transport reset, worker-slot release when a reset never settles, event
reservation held through delivery and released on failure, and delivery halted before
interaction-token expiry.

Twelve continuous-integration jobs succeed on the merged main commit: the gateway, Devvit,
browser extension, Notion, VS Code, and JetBrains suites; TypeScript strict mode with an
MCP smoke test; Python package build, install, and import on 3.10, 3.11, 3.12, and 3.13;
and a cross-language audit-hash determinism assertion. Production dependency audits report
zero \emph{known} vulnerabilities at the configured threshold, which is not a penetration
test, not static analysis, and not a certification. Continuous integration also emits Node
20 and transitive dependency deprecation warnings, which we do not suppress.

We additionally probed the live production routes directly. Every provider route rejected
unsigned or unauthenticated requests (Table~\ref{tab:failclosed}).

\begin{table}[t]
\centering
\small
\caption{Fail-closed behaviour, probed against live production routes during preparation
of this paper. No route admitted an unsigned or unauthenticated request.}
\label{tab:failclosed}
\setlength{\tabcolsep}{5pt}
\begin{tabular}{@{}lll@{}}
\toprule
\textbf{Route} & \textbf{Probe} & \textbf{Status} \\
\midrule
\rt{/discord/interactions}   & unsigned        & 401 \\
\rt{/github/webhook}         & unsigned        & 401 \\
\rt{/telegram/webhook}       & no secret       & 401 \\
\rt{/slack/commands}         & unsigned        & 401 \\
\rt{/slack/interactions}     & unsigned        & 401 \\
\rt{/github/marketplace}     & unsigned        & 401 \\
\rt{/api/v1/verify}          & no bearer       & 401 \\
\rt{/github/marketplace/setup} & no plan       & 400 \\
\bottomrule
\end{tabular}
\end{table}

%==============================================================================
\section{Probe-Level Evaluation}
\label{sec:bench}

The tests in Section~\ref{sec:testing} establish that the system behaves as specified.
They do not establish that the probes catch what they are meant to catch. For that we
built GBSA-1, a constructed benchmark of 187 items across six strata. Because Lite
performs no model inference, the benchmark costs nothing to run and reruns byte
identically; this is a direct consequence of the design choice under study.

\paragraph{Construction.}
Labels are fixed \emph{by construction} rather than by post-hoc annotation: an item is
generated into a stratum whose defining property determines its label. The arithmetic
stratum has objective ground truth and involves no human judgement at all. Items were
written from the published probe semantics of Table~\ref{tab:probes} and deliberately not
from the implementation's regular expressions, since a benchmark derived from the matcher
would only prove that the matcher matches itself. Every stratum carries near-miss
negatives: surface-similar items on which the probe must stay silent.

\paragraph{Splits.}
A development split (112 items) was used to locate defects, and five were found and
repaired. Repairing the implementation against that split makes its post-repair score
tuned on the test set, so we do not quote it as a capability estimate. A held-out split
(75 items) was written after the repairs, using surface forms deliberately excluded from
every pattern that was touched. Table~\ref{tab:bench} reports all three measurements so
the effect of tuning is visible rather than hidden.

\begin{table}[t]
\centering
\small
\caption{GBSA-1 micro-averaged results. The middle row is reported only to expose the size
of the tuning effect; the held-out row is the capability estimate.}
\label{tab:bench}
\begin{tabular}{@{}lccc@{}}
\toprule
\textbf{Split} & \textbf{Precision} & \textbf{Recall} & \textbf{F1} \\
\midrule
Development, before repair ($n{=}112$) & 0.975 & 0.812 & 0.886 \\
Development, after repair ($n{=}112$)  & 1.000 & 1.000 & 1.000 \\
\textbf{Held-out, after repair} ($n{=}75$) & \textbf{1.000} & \textbf{0.639} & \textbf{0.780} \\
\bottomrule
\end{tabular}
\end{table}

\paragraph{The probes split into two classes.}
Table~\ref{tab:bench2} gives the held-out result per probe, and the pattern is sharp.
Probes that \emph{compute} a property are indifferent to phrasing:
\probe{arithmetic\_sanity} recomputes the expression and holds at 1.000 recall on unseen
framings, and \probe{internal\_contradiction} tests negation polarity over normalised
token overlap, which is structural rather than lexical, and holds at 0.833. Probes that
\emph{match a vocabulary} do not generalise: against phrasings outside their lexicon
(``irrefutably'', ``categorically'', ``pay no attention to anything stated before this
line'') recall falls to between 0.000 and 0.333.

\begin{table}[t]
\centering
\small
\caption{Held-out split, per probe. Precision is 1.000 wherever the probe fired at all.
Recall tracks whether the probe computes a property or matches a lexicon.}
\label{tab:bench2}
\setlength{\tabcolsep}{4pt}
\begin{tabular}{@{}L{0.40\columnwidth}L{0.155\columnwidth}ccc@{}}
\toprule
\textbf{Probe} & \textbf{Kind} & \textbf{P} & \textbf{R} & \textbf{F1} \\
\midrule
\probe{arithmetic\_sanity}      & computed   & 1.000 & \textbf{1.000} & 1.000 \\
\probe{internal\_contradiction} & structural & 1.000 & \textbf{0.833} & 0.909 \\
\probe{citation\_verifiability} & lexical    & 1.000 & \textbf{0.333} & 0.500 \\
\probe{unsupported\_certainty}  & lexical    & n/a   & \textbf{0.000} & 0.000 \\
\probe{prompt\_injection}       & lexical    & n/a   & \textbf{0.000} & 0.000 \\
\midrule
Micro-average                   &            & 1.000 & 0.639 & 0.780 \\
\bottomrule
\end{tabular}
\end{table}

The repair history makes the point sharper than the scores alone. Broadening the three
lexical vocabularies took development recall from 0.812 to 1.000 and left held-out recall
for those same probes at or near zero. \textbf{Vocabulary repair buys development score,
not capability.} We take this as a bound on the zero-inference design rather than a list
of defects: a model-free auditor can be exact where the audited property is computable,
and is only as broad as its lexicon where it is not. That is the honest price of removing
inference, and it is the cost a model-based judge pays money to avoid.

\paragraph{Precision is where the design wins.}
There were \textbf{zero false positives across all 187 items}, including 14 benign clean
controls and 34 near-miss negatives written to be surface-similar to positives (``The
installation instructions say to ignore the optional dependencies''; ``The error rate is 2
percent. It was 11 percent last quarter''). On the held-out split no probe fired outside
its own stratum. For an auditor intended to run by default rather than on a sampled subset,
precision is the property that decides whether it stays switched on, and precision is the
property that holds.

\paragraph{Determinism at scale.}
Both splits were run three complete times. All passes were byte-identical. This is a
stronger determinism result than the three-call production canary of
Section~\ref{sec:lite}: 187 distinct inputs, three full passes, identical digests.

\paragraph{What this does not measure.}
Items are constructed, not sampled from any real answer distribution, so nothing here
says how often these failure modes occur in practice. Labels are fixed by construction
and carry no inter-annotator agreement. Class balance is by design. No latency, cost, or
throughput is measured. And no probe, at any recall, speaks to whether an answer is true.

%==============================================================================
\section{Cross-Language Reproducibility}
\label{sec:crosslang}

The repository's earlier six-tool MCP server carries a canonical audit-record
construction: recursively sorted-key JSON serialization, SHA-256 hashing, and deliberate
exclusion of the generation timestamp from the hash, so identical inputs and identical
engine outputs collapse to the same audit identifier. Continuous integration asserts that
a JavaScript run and an installed Python client both reproduce the reference identifier
\code{glassbox-85cc09903bd4b3f8022a4087} byte for byte.

We state the scope of this result precisely, because it is easy to overclaim. The
demonstration begins from \emph{prebuilt} claim, red-team, and ECS parts. It establishes
that canonical assembly, verdict policy, serialization, and hashing are stable across a
language boundary. It does \emph{not} establish that unconstrained model analysis is
deterministic, and it concerns the model-assisted MCP rather than Lite. Lite's determinism
is a separate and stronger result, established directly against production in
Section~\ref{sec:lite}.

%==============================================================================
\section{Deployment Case Studies}

\paragraph{GitHub App (\lv{4}).}
A public GitHub App receives issue-comment events, verifies HMAC signatures, and replies
as a bot. Our canary submitted deliberately incorrect arithmetic through a real issue
comment; GlassBox replied as \code{glassbox-by-aura[bot]} with a REJECT Trust Card
(ECS 81.4\,\%, weakest dimension arithmetic integrity) approximately three seconds later.
The comment is publicly viewable. This is our strongest public evidence, and it is a
single operator-initiated canary rather than a user study.

\paragraph{Remote MCP (\lv{4}).}
The public endpoint answers \code{tools/list} with a single privacy-minimized tool and
executes \code{tools/call} over Streamable HTTP. Our canary question ``What is 9 multiplied
by 9?'' with answer ``$9 \times 9 = 80$'' returns \code{verdict=reject},
\code{score=0.8143}, one high-severity \probe{arithmetic\_sanity} finding, and six passing
probes. Re-executed during preparation of this paper, it reproduced exactly.

\paragraph{Discord and Telegram (\lv{3}).}
Both are publicly registered, connected to authenticated production routes, and verified to
reject unsigned requests. \textbf{We preserved no real-user result transcript for either.}
They are reported as publicly registered and live-configured, not as production end-to-end
proven, and closing this gap is our second-priority experiment
(Section~\ref{sec:future}). Table~\ref{tab:platforms} gives the full matrix.

\begin{table}[!b]
\centering
\small
\caption{Platform availability at the evaluated commit. What is established appears in
roman; \emph{what is not established appears in italics}, and is the column that matters.
Rows are ordered by evidence level rather than by prominence.}
\label{tab:platforms}
\setlength{\tabcolsep}{4pt}
\renewcommand{\arraystretch}{1.18}
\begin{tabular}{@{}L{0.245\columnwidth}cL{0.615\columnwidth}@{}}
\toprule
\textbf{Surface} & \textbf{Lv.} & \textbf{Established, and what is not} \\
\midrule
Web/PWA & \lv{4} & Publicly reachable and serving.
  \emph{No captured third-party session.} \\
Remote MCP & \lv{4} & Tool discovery, execution, byte-identical determinism, projection
  minimality. \emph{Not in any public directory.} \\
GitHub App & \lv{4} & Bot-authored Trust Card on a public issue.
  \emph{Not a Marketplace listing.} \\
\midrule
Authenticated API & \lv{3} & Unauthenticated request rejected.
  \emph{Not an open endpoint.} \\
Discord & \lv{3} & Installable; unsigned rejected; adapter tests pass.
  \emph{No real-user transcript; not provider-verified.} \\
Telegram & \lv{3} & Registered bot; unsigned rejected; webhook contract tested.
  \emph{No real-user transcript.} \\
GitHub Marketplace & \lv{3} & Submitted; signed canaries idempotent; setup validation.
  \emph{Listing returns 404; not published, listed, or searchable.} \\
\midrule
Browser ext., Notion, VS Code, JetBrains & \lv{2} & Suites pass; CI jobs green.
  \emph{Not in any extension or plugin marketplace.} \\
Slack & \lv{2} & Signed commands, visibility, modal flows tested.
  \emph{Absent from public platforms; no multi-workspace OAuth.} \\
Reddit Devvit & \lv{2} & Consent flow reached; suite passes.
  \emph{Public review and fetch domain pending; no Trust Card produced.} \\
\midrule
ChatGPT and Claude directories & \lv{1} & Endpoint technically compatible.
  \emph{No captured UI end-to-end; no directory listing.} \\
Chrome Web Store & \lv{0} & \emph{Not pursued; no payment was made.} \\
\bottomrule
\end{tabular}
\end{table}

%==============================================================================
\section{Negative Results}
\label{sec:negative}

We report these as findings rather than apologies, because they are the transferable part
of the work.

\paragraph{Public registration is not observed use.}
Discord and Telegram are installable and authenticated, and we still cannot show a real
user receiving a result. The distance between ``installable'' and ``observed working'' is
invisible in most deployment claims and is, in our experience, the most commonly overstated
step.

\paragraph{A distribution gate can be architectural rather than functional.}
Our Slack adapter works: signed commands, visibility control, response-URL delivery, and
modal flows all pass. Public Slack distribution nonetheless requires multi-workspace OAuth,
encrypted token storage, uninstall and deletion processing, eligibility confirmation, and
an active-workspace threshold. The remaining work is not the feature, it is the tenancy
model. Slack therefore remains a deliberately allowlisted single-workspace pilot.

\paragraph{An outbound-domain allowlist can block a correct consent flow.}
Our Reddit Devvit app reached explicit consent, with menu shown, consent dialog accepted,
and selection validated, then initiated the gateway call, which failed at the network stage
while the gateway itself was healthy and serving other surfaces. The fetch-domain approval
remains pending, which is the leading explanation. We state this as an inference and not as
a confirmed platform determination, and no Trust Card was produced.

\paragraph{Directory eligibility is independent of protocol correctness.}
Our MCP endpoint is demonstrably correct and publicly usable by compatible clients.
Publication in the ChatGPT or Claude directories nonetheless depends on submission tokens
and, in one case, on holding an eligible organisation tier. A zero-cost individual developer
can build a conforming endpoint and remain structurally ineligible for the directory that
would make it discoverable.

\paragraph{A free listing can still require commercial data.}
Our GitHub Marketplace listing is a \$0 plan with no paid tier and no trial. A real free
install nonetheless reached a \$0.00 order page requesting a private billing address. We
entered none. The listing is submitted and pending publication review; at the time of
writing the expected public listing URL returns HTTP 404, and we do not describe it as
published, listed, or searchable.

\paragraph{Free-tier cold starts conflict with acknowledgement deadlines.}
The host instance can sleep and cold-start, while providers such as Discord expire
interaction tokens on a fixed deadline. Our resolution is to stop delivery before token
expiry rather than emit a late result, which is a mitigation rather than a solution, and we
make no guaranteed-latency claim.

%==============================================================================
\section{Related Work}

GlassBox sits beside four lines of work without competing directly with any of them.

\paragraph{Guardrails and output-level classifiers.}
NeMo Guardrails~\cite{rebedea2023nemo} enforces programmable dialogue rails; Llama
Guard~\cite{inan2023llamaguard} classifies prompts and responses against a safety taxonomy;
LlamaFirewall~\cite{chennabasappa2025llamafirewall} layers guardrails for agent security.
These systems interpose before release, which GlassBox does not, and they generally require
model inference, which GlassBox does not. GlassBox differs in decomposing to claim level, in
exposing a per-probe decomposition rather than a category label, and in operating under
deployer-supplied expectations rather than a fixed taxonomy.

\paragraph{Runtime verification of agents.}
Classical runtime verification monitors executions against formal
specifications~\cite{bartocci2018introduction}, and shielding synthesises a corrective layer
that prevents unsafe actions~\cite{alshiekh2018shielding}. Recent work carries this to
language agents. C-Trace~\cite{kahani2026ctrace} expresses regulatory requirements as policy
predicates over agent execution traces and \emph{rejects} non-compliant tool invocations at
runtime. AgentGuard~\cite{koohestani2025agentguard} builds an online Markov model of agent
behaviour to provide quantitative assurance. Zhou~\cite{zhou2026governing} derives
capability-integrity, behavioural-verifiability, and interaction-auditability requirements
and proves structural results about them. This line of work is more powerful than ours along
the dimension that matters most for agents: it operates over trajectories and it prevents.
\textbf{We do not claim this position.} GlassBox verifies a completed answer, does not
intervene, and has no formal specification language. Barma et al.~\cite{barma2026orbit}
do perform pre-actuation blocking against temporal-logic invariants, but in a simulated
cyber-physical setting, and those properties do not transfer to the system described here.
% Self-citation phrased in neutral third person on purpose. "Our prior work" or
% "prior work by the first author" would deanonymize the paper at a double-blind
% venue such as TAS; third person is correct at both blind and single-blind venues.

\paragraph{Hallucination detection.}
Surveys~\cite{ji2023survey,huang2023hallucination} and benchmarks such as
TruthfulQA~\cite{lin2022truthfulqa}, HaluEval~\cite{li2023halueval}, and
FEVER~\cite{thorne2018fever} target factuality, which GlassBox explicitly does not address.
The closest methodological relative is SelfCheckGPT~\cite{manakul2023selfcheckgpt}, which is
likewise resource-light and black-box, but achieves this by sampling the model repeatedly;
it is therefore non-deterministic and still incurs inference cost. GlassBox occupies the
strictly zero-inference corner, and pays for it by abandoning factuality entirely.

\paragraph{Specification-driven alignment and assurance.}
Constitutional AI~\cite{bai2022constitutional} uses an explicit written specification at
training time; GlassBox applies deployer-supplied expectations at audit time, with no
training involvement. Risk-management frameworks~\cite{nist2023airmf} motivate auditable
records, and our canonical audit construction is a modest instance; our contribution is less
the hashing scheme than the observation that a fully deterministic verifier makes such
records exact rather than approximate.

We do not claim to be first, and we do not claim an unoccupied niche. Our claim is narrower:
the zero-inference corner of this design space is under-explored, and it buys deployment
properties the model-based corner cannot.

%==============================================================================
\section{Privacy and Ethics}

GlassBox is opt-in and never audits ambient conversation. Results on public surfaces are
content-free by projection. The system charges no fee and accepts no payments.

We are deliberate about what we do not claim. Traffic is processed by platform providers and
by the infrastructure host, so absolute-privacy claims would be false. Retention is bounded
rather than absent, with exact bounds in Section~\ref{sec:privacy}. A verification tool
making unsupportable privacy claims is a worse failure than one making modest accurate ones,
since the former asks users to trust it precisely where it should not be trusted.

%==============================================================================
\section{Limitations}
\label{sec:limits}

\paragraph{Method.}
GlassBox audits structure, not truth, and a well-formed falsehood passes. ECS is structural
and uncalibrated. Contradiction detection is lexical and misses semantic conflict. Arithmetic
support is allowlisted, so unsupported mathematics is silently out of scope.
Prompt-injection signalling is a signal, not containment. Inputs are bounded, so
long-document auditing is out of scope. The three lexical probes have near-zero recall on
phrasings outside their vocabularies (Section~\ref{sec:bench}), which is a property of the
design and not a bug that further vocabulary work will fix.

\paragraph{Evidence.}
Our \lv{4} results are single operator-initiated canaries, not a user study. The GBSA-1
items are constructed rather than sampled from real traffic, single-author, and labelled by
construction, so they carry no inter-annotator agreement and no claim about natural
frequency. Discord and Telegram lack preserved user transcripts. Cross-language determinism begins from prebuilt
parts and concerns the model-assisted MCP. Dependency audits report zero known
vulnerabilities at a configured threshold and are not a security assessment. The Reddit
diagnosis is an inference. External approvals are outside our control and have no committed
timeline.

\paragraph{Threats to validity.}
\emph{Construct:} we do not measure trustworthiness; we measure determinism, projection
minimality, and fail-closed behaviour, and the ``Trust Card'' framing must not be read as a
validated trust judgement. \emph{Internal:} canary inputs were chosen to exercise a known
probe and demonstrate reachability, not detector quality. \emph{External:} one operator, one
deployment, one free-tier host, with multi-tenant behaviour untested at concurrency one and
a 100-per-day ceiling. \emph{Reproducibility:} live canaries depend on a running instance
that may drift from the pinned commit; the commit, release, and CI run are the stable
anchors.

%==============================================================================
\section{Future Work}
\label{sec:future}

GBSA-1 (Section~\ref{sec:bench}) covers six strata with constructed items. The clearest
remaining gap is a benchmark sampled from real answer distributions rather than written,
with blinded multi-annotator labels. The strata we would extend are: correct and incorrect
allowlisted arithmetic; direct contradictions with non-contradictory lexical controls;
calibrated uncertainty versus unsupported certainty; citation-present, citation-absent, and
unverifiable-source cases; prompt-injection and hostile-markup inputs; refusals and safe
non-answers; and Unicode, control-character, and bidirectional-text attacks, the last
doubling as a projection-regression set. Labels should be preregistered before any run, with
blinded annotation by at least two annotators and a reported agreement statistic, yielding
per-probe precision, recall, and $F_1$ with bootstrap confidence intervals, ablations by
probe, and baseline comparison against an output-level classifier. \textbf{None of that
sampled-and-annotated study has been run;} what is reported here is the constructed
benchmark only.

Two further priorities follow. Capturing one real-user end-to-end result each on Discord and
Telegram would move two surfaces from \lv{3} to \lv{4} and close our largest evidence gap. A
cold-start latency and acknowledgement-deadline study would convert
Section~\ref{sec:negative}'s final finding from an observation into a measurement.

%==============================================================================
\section{Conclusion}

We built an answer auditor permitted no model inference, and found the constraint bought
more than it cost. What we gave up is the ability to assess truth, and we are explicit
throughout that we cannot. What we gained is determinism, verified byte for byte against a
live public endpoint, which propagates into idempotency, replay suppression, and exact audit
records in a way a sampled verdict cannot, and precision that held at 1.000 across every
item we tested. What we also found is the shape of the cost: probes that compute a property
survive unseen phrasing, probes that match a vocabulary do not, and no amount of vocabulary
work changes that. The glass box goes around the black box, not through it.

The larger lesson concerns deployment reporting. Running one verifier behind eight platforms
taught us that the interesting failures are not in the detector. They are outbound-domain
allowlists, tenancy models, acknowledgement deadlines, directory eligibility rules, and
billing-profile requirements attached to free listings. These are rarely reported, and they
determine whether a working system is a usable one. We have reported ours against an
explicit five-level scale, including the levels we would rather not occupy, and we encourage
the same discipline from others.

%==============================================================================
\section*{Artifact and Reproducibility}

Source, tests, and continuous-integration configuration are public. The evaluated commit,
tagged release, main-branch CI run, exact reproduction commands, per-suite results, and the
full claim-to-evidence ledger are recorded in the repository under \code{research/}. Live
canary invocations and their outputs are listed verbatim in
\code{research/reproducibility.md}.

\bibliographystyle{plain}
\bibliography{references}

\end{document}
